\begin{proposition}[2.6.2]
\label{IV.2.6.2-fr}
Con las notaciones de \textup{(2.6.1)}, supongamos que el morfismo $g:S'\to S$ es
fielmente plano y cuasicompacto. Consideremos, para un morfismo, la propiedad
de ser:
\begin{enumerate}
\item[\textup{(i)}] abierto;
\item[\textup{(ii)}] cerrado;
\item[\textup{(iii)}] cuasicompacto y un homeomorfismo sobre el subespacio imagen;
\item[\textup{(iv)}] un homeomorfismo.
\end{enumerate}
Entonces, si $\boldsymbol{P}$ designa una de las propiedades precedentes y si $f'$ posee la
propiedad $\boldsymbol{P}$, lo mismo ocurre con $f$.
\end{proposition}

\oldpage[IV-2]{28}
Como el morfismo $Y'\to Y$ es fielmente plano y cuasicompacto (2.2.13 y
1.1.2), es posible, en virtud de (\textbf I, 3.3.11), limitarse al caso en que $Y=S$,
$Y'=S'$. Si $g'$ es la proyección $X'\to X$, se sabe que, para toda parte $M$
de $X$, se tiene $g^{-1}(f(M))=f'(g'^{-1}(M))$ (\textbf{I}, 3.4.8). Si $f'$ es un morfismo abierto (resp. cerrado), entonces, para toda parte
abierta (resp. cerrada) $M$ de $X$, $f'(g'^{-1}(M))$ es abierta (resp.
cerrada) en $Y'$, y, como $g$ es fielmente plano y cuasicompacto, se concluye
que $f(M)$ es abierta (resp. cerrada) en $Y$ en virtud de (2.3.12). Esto
demuestra la proposición en los casos (i) y (ii). Demostrémosla en el caso
(iii) (lo que la demostrará en el caso (iv), teniendo en cuenta (2.6.1,
(iv)). En virtud de (2.6.1, (ii)), $f$ es inyectivo, y queda por demostrar que
$f$, considerado como aplicación de $X$ sobre $f(X)$, es una aplicación
cuasicompacta y \emph{abierta}. Como $f'$ es una aplicación cuasicompacta, lo mismo ocurre con $f$
(1.1.4). Basta entonces demostrar que, para toda parte cerrada $Z$ de $X$, se
tiene $Z=f^{-1}\overline{(f(Z))}$; como $g'$ es sobreyectivo, esta relación equivale a
$g'^{-1}(Z)=g'^{-1}(f^{-1}\overline{(f(Z))})$ o bien a
$g'^{-1}(Z)=f'^{-1}(g^{-1}\overline{(f(Z))})$. Ahora bien, como $f$ es cuasicompacto, lo mismo vale para su composición con la
inyección canónica $Z\to X$ (siendo aquí $Z$ el subpreesquema cerrado reducido
de $X$ que tiene a $Z$ como espacio subyacente). Aplicando (2.3.10) a la parte
$f(Z)$ de $Y$ (imagen del morfismo $f|Z$), se obtiene $g^{-1}\overline{(f(Z))}=\overline{g^{-1}(f(Z))}$;
la fórmula que hay que demostrar equivale, por tanto, a la siguiente:
$Z'=f'^{-1}\overline{(f'(Z'))}$, donde se ha puesto $Z'=g'^{-1}(Z)$; pero esta fórmula se deduce de la hipótesis
de que $f'$ es un homeomorfismo de $X'$ sobre $f'(X')$.
